\ifdefined\gkver\else\def\gkver{student}\fi
\documentclass[\gkver]{gaokaozhenti}
\begin{document}

\chapter{1988年全国}

\begin{enumerate}
\item
%% number: 1
%% paperName: 1988年全国高考第 1 题 2 分
%% typeId: 1
%% type: 单选题
%% chapter: 运动和力
%% point: 超重失重
%% method: 无
%% score: 2
%% degree: 850
%% duplicateId: 0
%% body:
电梯内有一个物体，质量为 $m$，用细线挂在电梯的天花板上。当电梯以 $\frac{g}{3}$ 的加速度竖直加速下降时（$g$ 为重力加速度），细线对物体的拉力为\xzanswer{A}
\fourchoices[answer=A]
{$\frac{2}{3}mg$}
{$\frac{1}{3}mg$}
{$\frac{4}{3}mg$}
{$mg$}

%% answer: A
%% memo:
\memoanswer{
由题意知，设绳力为 $T$，由牛顿第二定律得 $mg-T=ma$，又 $a=\frac{g}{3}$，解得 $T=\frac{2}{3}mg$，故 A 正确。
}

\item
%% number: 2
%% paperName: 1988年全国高考第 2 题 3 分
%% typeId: 1
%% type: 单选题
%% chapter: 功和能
%% point: 机械能守恒定律
%% method: 无
%% score: 3
%% degree: 850
%% duplicateId: 0
%% body:
一个人站在阳台上，以相同的速度 $v_{0}$ 分别把三个球竖直向上抛出、竖直向下抛出、水平抛出，不计空气阻力，则三球落地时的速率\xzanswer{D}
\fourchoices[answer=D]
{上抛球最大}
{下抛球最大}
{平抛球最大}
{三球一样大}

%% answer: D
%% memo:
\memoanswer{
设落地速率为 $v$，由机械能守恒得 $mgh=\frac{1}{2}mv^{2}$，解得 $v=\sqrt{2gh}$，故 D 正确。
}

\item
%% number: 3
%% paperName: 1988年全国高考第 3 题 2 分
%% typeId: 1
%% type: 单选题
%% chapter: 运动和力
%% point: 连接体
%% method: 无
%% score: 2
%% degree: 650
%% duplicateId: 0
%% body:
两物体 A 和 B，质量分别为 $m_{1}$ 和 $m_{2}$，互相接触放在光滑水平面上，如图所示。对物体 A 施以水平的推力 $F$，则物体 A 对物体 B 的作用力等于\xzanswer{B}
\onepicture{\includesvg{figs/03}}
\fourchoices[answer=B]
{$\frac{m_{1}}{m_{1}+m_{2}}F$}
{$\frac{m_{2}}{m_{1}+m_{2}}F$}
{$F$}
{$\frac{m_{2}}{m_{1}}F$}

%% answer: B
%% memo:
\memoanswer{
对系统整体分析，有 $F=(m_{1}+m_{2})a$；隔离 B 有 $F_{AB}=m_{2}a$，联立解得 $F_{AB}=\frac{m_{2}}{m_{1}+m_{2}}F$，故 B 正确。
}

\item
%% number: 4
%% paperName: 1988年全国高考第 4 题 2 分
%% typeId: 1
%% type: 单选题
%% chapter: 光学
%% point: 光的电磁说
%% method: 无
%% score: 2
%% degree: 850
%% duplicateId: 0
%% body:
单色光从真空射入玻璃时，它的\xzanswer{D}
\fourchoices[answer=D]
{波长变长，波速变小}
{波长变短，波速变大}
{波长变长，波速变大}
{波长变短，波速变小}

%% answer: D
%% memo:
\memoanswer{
由光的折射特性知 $v=\frac{c}{n}$，因 $n>1$，所以 $v<c$。又 $c=f\lambda_{0}$，$v=f\lambda_{1}$，而光折射是频率 $f$ 不变，所以 $\lambda_{1}<\lambda_{0}$，故 D 正确。
}

\item
%% number: 5
%% paperName: 1988年全国高考第 5 题 2 分
%% typeId: 1
%% type: 单选题
%% chapter: 原子和原子核
%% point: 原子核式结构
%% method: 无
%% score: 2
%% degree: 950
%% duplicateId: 0
%% body:
原子的核式结构学说，是卢瑟福根据以下哪个实验或现象提出来的？\xzanswer{C}
\fourchoices[answer=C]
{光电效应实验}
{氢原子光谱实验}
{$\alpha$ 粒子散射实验}
{天然放射现象}

%% answer: C
%% memo:
\memoanswer{
卢瑟福根据 $\alpha$ 粒子散射实验提出原子的核式结构学说，故 C 正确。
}

\item
%% number: 6
%% paperName: 1988年全国高考第 6 题 2 分
%% typeId: 1
%% type: 单选题
%% chapter: 直线运动
%% point: 竖直上下抛运动
%% method: 无
%% score: 2
%% degree: 850
%% duplicateId: 0
%% body:
将一物体以某一初速竖直上抛，在下列四幅图中，哪一幅能正确表示物体在整个运动过程中的速率 $v$ 与时间 $t$ 的关系（不计空气阻力）？\xzanswer{B}
\fourchoices[answer=B, ispicture=true, h=2.6cm]
{\includesvg{figs/06a}}
{\includesvg{figs/06b}}
{\includesvg{figs/06c}}
{\includesvg{figs/06d}}

%% answer: B
%% memo:
\memoanswer{
由小球运动的规律知 $v=v_{0}-gt$，所以速率先变小后变大，故 B 正确。
}

\item
%% number: 7
%% paperName: 1988年全国高考第 7 题 2 分
%% typeId: 1
%% type: 单选题
%% chapter: 电磁感应
%% point: 楞次定律推广
%% method: 无
%% score: 2
%% degree: 650
%% duplicateId: 0
%% body:
在水平放置的光滑绝缘杆 ab 上，挂有两个金属环 M 和 N，两环套在一个通电密绕长螺线管的中部，如图所示。螺线管中部区域的管外磁场可以忽略。当变阻器的滑动接头向左移动时，两环将怎样运动？\xzanswer{C}
\onepicture{\includesvg{figs/07}}
\fourchoices[answer=C]
{两环一起向左移动}
{两环一起向右移动}
{两环互相靠近}
{两环互相离开}

%% answer: C
%% memo:
\memoanswer{
接头向左滑动时，电阻 $R$ 减小，电流 $I$ 变大，管内磁通量向右增大，金属环中感生电流的磁场向左，两环感应电流都沿顺时针方向（从左向右看），所以两环互相靠近，故 C 正确。
}

\item
%% number: 8
%% paperName: 1988年全国高考第 8 题 2 分
%% typeId: 1
%% type: 单选题
%% chapter: 电路
%% point: 电阻定律
%% method: 无
%% score: 2
%% degree: 650
%% duplicateId: 0
%% body:
把两根同种材料的电阻丝分别连在两个电路中。甲电阻丝长为 $l$，直径为 $d$；乙电阻丝长为 $2l$，直径为 $2d$。要使两电阻丝消耗的功率相等。加在两电阻丝上的电压比应满足\xzanswer{C}
\fourchoices[answer=C]
{$U_{\text{甲}}:U_{\text{乙}}=1$}
{$U_{\text{甲}}:U_{\text{乙}}=\sqrt{2}:2$}
{$U_{\text{甲}}:U_{\text{乙}}=\sqrt{2}$}
{$U_{\text{甲}}:U_{\text{乙}}=2$}

%% answer: C
%% memo:
\memoanswer{
由电阻定律得 $R=\rho\frac{L}{S}$，又 $S=\frac{1}{4}\pi d^{2}$，得 $\frac{R_{\text{甲}}}{R_{\text{乙}}}=\frac{1}{2}\times\left(\frac{2}{1}\right)^{2}=\frac{2}{1}$，由功率表达式得 $P=\frac{U^{2}}{R}$，得 $U=\sqrt{PR}$，所以 $\frac{U_{\text{甲}}}{U_{\text{乙}}}=\sqrt{\frac{R_{\text{甲}}}{R_{\text{乙}}}}=\frac{\sqrt{2}}{1}$，故 C 正确。
}

\item
%% number: 9
%% paperName: 1988年全国高考第 9 题 2 分
%% typeId: 1
%% type: 单选题
%% chapter: 交流电
%% point: LC振荡回路
%% method: 无
%% score: 2
%% degree: 850
%% duplicateId: 0
%% body:
要使 LC 振荡电路的周期增大一倍。可采用的办法是\xzanswer{A}
\fourchoices[answer=A]
{自感系数 $L$ 和电容 $C$ 都增大一倍}
{自感系数 $L$ 和电容 $C$ 都减小一半}
{自感系数 $L$ 增大一倍，而电容 $C$ 减小一半}
{自感系数 $L$ 减小一半，而电容 $C$ 增大一倍}

%% answer: A
%% memo:
\memoanswer{
LC 振荡电路的周期公式为 $T=2\pi\sqrt{LC}$，当 $L$ 与 $C$ 都增大 1 倍时，则 $T$ 增大 1 倍，故 A 正确；当 $L$ 与 $C$ 都减小一半时，则 $T$ 减小一半，故 B 错误；当 $L$ 增大 1 倍，$C$ 减小一半时，则 $T$ 不变，故 C 错误；当 $L$ 减小一半，$C$ 增大 1 倍时，则 $T$ 不变，故 D 错误。
}

\item
%% number: 10
%% paperName: 1988年全国高考第 10 题 2 分
%% typeId: 1
%% type: 单选题
%% chapter: 力物体的平衡
%% point: 物体系平衡
%% method: 整体与隔离
%% score: 2
%% degree: 650
%% duplicateId: 0
%% body:
在粗糙水平面上有一个三角形木块 abc，在它的两个粗糙斜面上分别放两个质量 $m_{1}$ 和 $m_{2}$ 的木块，$m_{1}>m_{2}$，如图所示。已知三角形木块和两物体都是静止的，则粗糙水平面对三角形木块\xzanswer{D}
\onepicture{figs/10.png}
\fourchoices[answer=D]
{有摩擦力的作用，摩擦力的方向水平向右}
{有摩擦力的作用，摩擦力的方向水平向左}
{有摩擦力的作用，但摩擦力的方向不能确定，因为 $m_{1}$、$m_{2}$、$\theta_{1}$、$\theta_{2}$ 的数值并未给出}
{以上结论都不对}

%% answer: D
%% memo:
\memoanswer{
隔离 $m_{1}$，压力 $N_{1}=m_{1}g\cos\theta_{1}$，摩擦力 $f_{1}=m_{1}g\sin\theta_{1}$；同理，对 $m_{2}$，压力 $N_{2}=m_{2}g\cos\theta_{2}$，摩擦力为 $f_{2}=m_{2}g\sin\theta_{2}$；三角形木块在水平方向的受力为 $F_{x}=(f_{2}\cos\theta_{2}-N_{2}\sin\theta_{2})-(f_{1}\cos\theta_{1}-N_{1}\sin\theta_{1})=0$，所以无摩擦力作用，故 D 正确。
}

\item
%% number: 11
%% paperName: 1988年全国高考第 11 题 2 分
%% typeId: 1
%% type: 单选题
%% chapter: 曲线运动
%% point: 人造地球卫星
%% method: 无
%% score: 2
%% degree: 550
%% duplicateId: 0
%% body:
两个球形行星 A 和 B 各有一卫星 a 和 b，卫星的圆轨道接近各自行星的表面。如果两行星质量之比 $M_{A}/M_{B}=p$，两行星半径之比 $R_{A}/R_{B}=q$，则两卫星周期之比 $T_{a}/T_{b}$ 为\xzanswer{A}
\fourchoices[answer=A]
{$q\sqrt{\frac{q}{p}}$}
{$q\sqrt{p}$}
{$p\sqrt{\frac{p}{q}}$}
{$\sqrt{pq}$}

%% answer: A
%% memo:
\memoanswer{
万有引力提供向心力得 $G\frac{Mm}{R^{2}}=m\cdot\frac{4\pi^{2}}{T^{2}}R$，得 $T=2\pi\sqrt{\frac{R^{3}}{GM}}$，所以 $\frac{T_{A}}{T_{B}}=\sqrt{\left(\frac{R_{A}}{R_{B}}\right)^{3}\cdot\frac{M_{B}}{M_{A}}}=q\sqrt{\frac{q}{p}}$，故 A 正确。
}

\item
%% number: 12
%% paperName: 1988年全国高考第 12 题 2 分
%% typeId: 1
%% type: 单选题
%% chapter: 交流电
%% point: 变压器
%% method: 无
%% score: 2
%% degree: 650
%% duplicateId: 0
%% body:
图中 S 为三相交流电源。连在电路中的各电阻丝 B$_{1}$、B$_{2}$、B$_{3}$、B$_{4}$ 的阻值都相同。变压器都是电压比为 2$:$1 的降压变压器。这些电阻丝，按消耗功率大小的顺序排列（从大到小）应为\xzanswer{B}
\onepicture{figs/12.png}
\fourchoices[answer=B]
{B$_{1}$B$_{2}$B$_{3}$B$_{4}$}
{B$_{4}$B$_{3}$B$_{2}$B$_{1}$}
{B$_{4}$B$_{2}$B$_{3}$B$_{1}$}
{B$_{3}$B$_{1}$B$_{4}$B$_{2}$}

%% answer: B
%% memo:
\memoanswer{
变压器初级电压为 $U_{1}=U_{\text{相}}$，$U_{2}=U_{\text{线}}$，变压器次级电压为 $U_{1}'=\frac{U_{\text{相}}}{2}$，$U_{2}'=\frac{U_{\text{线}}}{2}$，又 $U_{3}=U_{\text{相}}$，$U_{4}=U_{\text{线}}$，线电压与相电压的关系 $U_{\text{线}}=\sqrt{3}U_{\text{相}}$，又功率 $P=\frac{U^{2}}{R}$，所以 $P_{1}:P_{2}:P_{3}:P_{4}=\frac{1}{4}:\frac{3}{4}:1:3$，所以这些电阻丝按消耗功率从大到小的顺序排列应为 $B_{4}B_{3}B_{2}B_{1}$，故 B 正确。
}

\item
%% number: 13
%% paperName: 1988年全国高考第 13 题 3 分
%% typeId: 2
%% type: 多选题
%% chapter: 原子和原子核
%% point: 人工核反应
%% method: 无
%% score: 3
%% degree: 850
%% duplicateId: 0
%% body:
下列核反应方程中，哪些是平衡的？\xzanswer{BC}
\fourchoices[answer=BC]
{$\ce{^{10}_{5}B}$ $+$（$\alpha$ 粒子）$\to$ $\ce{^{13}_{7}N}$ $+$ $\ce{^{1}_{1}H}$}
{$\ce{^{27}_{13}Al}$ $+$ $\ce{^{1}_{0}n}$ $\to$ $\ce{^{27}_{12}Ma}$ $+$（中子）}
{$\ce{^{9}_{4}Be}$ $+$ $\ce{^{4}_{2}He}$ $\to$ $\ce{^{12}_{6}C}$ $+$（中子）}
{$\ce{^{14}_{7}N}$ $+$ $\ce{^{4}_{2}He}$ $\to$ $\ce{^{17}_{8}O}$ $+$（氘核）}

%% answer: BC
%% memo:
\memoanswer{
核反应方程满足质量数守恒和核电荷数守恒，分析每个选项后知 BC 正确。
}

\item
%% number: 14
%% paperName: 1988年全国高考第 14 题 3 分
%% typeId: 2
%% type: 多选题
%% chapter: 气体性质
%% point: 分子动理论
%% method: 无
%% score: 3
%% degree: 850
%% duplicateId: 0
%% body:
在有关布朗运动的说法中，正确的是\xzanswer{BC}
\fourchoices[answer=BC]
{液体的温度越低，布朗运动越显著}
{液体的温度越高，布朗运动越显著}
{悬浮微粒越小，布朗运动越显著}
{悬浮微粒越大，布朗运动越显著}

%% answer: BC
%% memo:
\memoanswer{
液体的温度越高，液体分子运动越剧烈，布朗运动越显著，故 A 错误 B 正确；悬浮微粒越小，同时液体分子碰撞的几率小，不平衡性越显著，布朗运动越显著，故 C 正确 D 错误。
}

\item
%% number: 15
%% paperName: 1988年全国高考第 15 题 3 分
%% typeId: 2
%% type: 多选题
%% chapter: 光学
%% point: 光的干涉
%% method: 无
%% score: 3
%% degree: 850
%% duplicateId: 0
%% body:
下列哪些现象说明光具有波动性？\xzanswer{AB}
\fourchoices[answer=AB]
{光的干涉}
{光的衍射}
{光的反射}
{光电效应}

%% answer: AB
%% memo:
\memoanswer{
光的干涉和光的衍射现象说明光具有波动性，故 AB 正确；光电效应说明光具有粒子性，故 D 错误；光的反射不能说明光具有波动性，故 C 错误。
}

\item
%% number: 16
%% paperName: 1988年全国高考第 16 题 3 分
%% typeId: 1
%% type: 单选题
%% chapter: 光学
%% point: 透镜
%% method: 无
%% score: 3
%% degree: 650
%% duplicateId: 0
%% body:
如图表示的是透镜成像实验的装置。\xzanswer{D}
\onepicture{figs/16.png}
\fourchoices[answer=D]
{如透镜是凸透镜，则不论物体放在透镜左方何处，只要把光屏移到适当位置，一定能在屏上得到物体的像}
{如透镜是凸透镜，则不论物体放在透镜左方何处，去掉光屏而用眼睛从右向左沿主轴直接观察，一定看不到物体的像}
{如透镜是凹透镜，则不论物体放在透镜左方何处，只要把光屏移到适当位置，一定能在屏上得到物体的像}
{如透镜是凹透镜，则不论物体放在透镜左方何处，去掉光屏而用眼睛从右向左沿主轴直接观察，一定能看到物体的像}

%% answer: D
%% memo:
\memoanswer{
当物体放在透镜焦距之内时，屏上得不到物体的像，故 A 错误；当物体放在透镜焦距之外时，用眼睛看不到物体的像，故 B 错误；物体放在凹透镜左侧时，一定不能在屏上得到物体的像，故 C 错误；物体放在凹透镜左侧时，成与物同侧的虚像，所以能看到，故 D 正确。
}

\item
%% number: 17
%% paperName: 1988年全国高考第 17 题 3 分
%% typeId: 1
%% type: 单选题
%% chapter: 电路
%% point: 电容
%% method: 无
%% score: 3
%% degree: 650
%% duplicateId: 0
%% body:
两块平行金属板带等量异号电荷，要使两板间的电压加倍，而板间的电场强度减半，采用的办法有\xzanswer{C}
\fourchoices[answer=C]
{两板的电量加倍，而距离变为原来的 4 倍}
{两板的电量加倍，而距离变为原来的 2 倍}
{两板的电量减半，而距离变为原来的 4 倍}
{两板的电量减半，而距离变为原来的 2 倍}

%% answer: C
%% memo:
\memoanswer{
平行板电容器决定式为 $C=\frac{\varepsilon S}{4\pi kd}$，又 $U=\frac{Q}{C}$，$E=\frac{U}{d}$，故 $E=\frac{Q}{Cd}=\frac{4\pi k}{\varepsilon}\cdot\frac{Q}{S}$，则 $E$ 与板间距 $d$ 无关，仅与 $Q$、$S$、$\varepsilon$ 有关。当 $Q$ 加倍时，$E$ 加倍，故 AB 错误；当 $Q$ 减半时，$E$ 减小为原来的一半，故 C 正确 D 错误。
}

\item
%% number: 18
%% paperName: 1988年全国高考第 18 题 3 分
%% typeId: 2
%% type: 多选题
%% chapter: 电磁感应
%% point: 线圈过有界磁场
%% method: 无
%% score: 3
%% degree: 750
%% duplicateId: 0
%% body:
如图所示，闭合矩形线圈 abcd 从静止开始竖直下落，穿过一个匀强磁场区域，此磁场区域竖直方向的长度远大于矩形线圈 bc 边的长度。不计空气阻力，则\xzanswer{BC}
\onepicture{\includesvg{figs/18}}
\fourchoices[answer=BC]
{从线圈 dc 边进入磁场到 ab 边穿出磁场的整个过程，线圈中始终有感应电流}
{从线圈 dc 边进入磁场到 ab 边穿出磁场的整个过程中，有一个阶段线圈的加速度等于重力加速度}
{dc 边刚进入磁场时线圈内感应电流的方向，与 dc 边刚穿出磁场时感应电流的方向相反}
{dc 边刚进入磁场时线圈内感应电流的大小，与 dc 边刚穿出磁场时感应电流的大小一定相等}

%% answer: BC
%% memo:
\memoanswer{
当线圈全部在磁场中时，磁通量不变，无感应电流，故 A 错误；当线圈全部在磁场中时，磁通量不变，无感应电流，无安培力，$a=g$，故 B 正确；dc 边进入磁场时，产生感生电动势，感应电流逆时针，ab 穿出磁场时，产生感应电动势，电流感应顺时针，方向相反，故 C 正确；感生电流 $I=\frac{BLv}{R}$，因速度不等，所以感生电流大小不等，故 D 错误。
}

\item
%% number: 19
%% paperName: 1988年全国高考第 19 题 3 分
%% typeId: 2
%% type: 多选题
%% chapter: 磁场
%% point: 带电粒子在复合场中的运动
%% method: 无
%% score: 3
%% degree: 650
%% duplicateId: 0
%% body:
设空间存在竖直向下的匀强电场和垂直纸面向里的匀强磁场，如图所示。已知一离子在电场力和洛仑兹力的作用下，从静止开始自 A 点沿曲线 ACB 运动，到达 B 点时速度为零。C 点是运动的最低点。忽略重力，以下说法中正确的是\xzanswer{ABC}
\onepicture{figs/19.png}
\fourchoices[answer=ABC]
{这离子必带正电荷}
{A 点和 B 点位于同一高度}
{离子在 C 点时速度最大}
{离子到达 B 点后，将沿原曲线返回 A 点}

%% answer: ABC
%% memo:
\memoanswer{
根据粒子的轨迹，由左手定则判断离子带正电荷，故 A 正确；因洛伦兹力不做功，A、B 重力势能相等，所以在同一高度，故 B 正确；重力势能转化为动能，轨迹中 C 点的重力势能最小，故动能最大，即速度最大，故 C 正确；由左手定则知，离子到达 B 点后向右做相似轨迹的运动（摆线运动），故 D 错误。
}

\item
%% number: 20
%% paperName: 1988年全国高考第 20 题 3 分
%% typeId: 2
%% type: 多选题
%% chapter: 运动和力
%% point: 加速减速模型
%% method: 无
%% score: 3
%% degree: 550
%% duplicateId: 0
%% body:
一物体放在光滑水平面上，初速为零。先对物体施加一向东的恒力 $F$，历时 1 秒钟；随即把此力改为向西，大小不变，历时 1 秒钟；接着又把此力改为向东，大小不变，历时 1 秒钟；如此反复，只改变力的方向，共历时 1 分钟。在此 1 分钟内\xzanswer{D}
\fourchoices[answer=D]
{物体时而向东运动，时而向西运动，在 1 分钟末静止于初始位置之东}
{物体时而向东运动，时而向西运动，在 1 分钟末静止于初始位置}
{物体时而向东运动，时而向西运动，在 1 分钟末继续向东运动}
{物体一直向东运动，从不向西运动，在 1 分钟末静止于初始位置之东}

%% answer: D
%% memo:
\memoanswer{
物体在第 1 个 1 s 向东做匀加速运动，在第 2 个 1 s 向东做匀减速运动，速度恰好减至 0，后续运动与前 2 s 的运动规律相同，所以物体一直向东运动，从不向西运动，在 1 分钟末静止于初始位置之东，故 D 正确。
}

\item
%% number: 21
%% paperName: 1988年全国高考第 21 题 3 分
%% typeId: 3
%% type: 填空题
%% chapter: 电路
%% point: 串并联电路
%% method: 无
%% score: 3
%% degree: 850
%% duplicateId: 0
%% body:
如图所示电路中，$R_{0}$ 是已知的，要使 AB 间的总电阻恰等于 $R_{0}$，则 $R_{1}=$\tkanswer[1.4cm]{$\frac{\sqrt{3}}{3}R_{0}$}。
\onepicture[align=right]{figs/21.png}

%% answer: \frac{\sqrt{3}}{3}R_{0}
\jdanswer{
$\frac{\sqrt{3}}{3}R_{0}$
}

%% memo:
\memoanswer{
由题意知 $R_{AB}=R_{1}+\frac{R_{1}(R_{1}+R_{0})}{2R_{1}+R_{0}}=\frac{3R_{1}^{2}+2R_{1}R_{0}}{2R_{1}+R_{0}}$，令 $R_{AB}=R_{0}$，解得 $R_{1}=\frac{\sqrt{3}}{3}R_{0}$。
}

\item
%% number: 22
%% paperName: 1988年全国高考第 22 题 3 分
%% typeId: 3
%% type: 填空题
%% chapter: 气体性质
%% point: 阿伏伽德罗常数
%% method: 无
%% score: 3
%% degree: 550
%% duplicateId: 0
%% body:
有一真空容器，在室温下容器内的气压为 $10^{-8}\UPa$。估算该容器内 $1\Ucmc$ 气体中的分子数。估算取 1 位有效数字即可。答：\tkanswer[2cm]{$3\times10^{6}$}。

1 标准大气压 $=1\times10^{5}\UPa$。阿伏伽德罗常数 $N=6\times10^{23}\Umoln$。

%% answer: 3×10^6
\jdanswer{
$3\times10^{6}$
}

%% memo:
\memoanswer{
标准状况下，1 mol 气体的体积为 $22.4\UL=22.4\times10^{-3}\Umc$。

$1\Ucmc$ 该气体室温下的压强 $p_{1}=10^{-8}\UPa$，体积为 $V_{1}=1\Ucmc=1\times10^{-6}\Umc$，室温取 $0\UdC$，$T_{1}=273\UK$。

$1\Ucmc$ 该气体在标准状况下的体积为 $V_{2}$，压强为 $p_{2}=1\times10^{5}\UPa$，温度为 $T_{2}=273\UK$，由玻意耳定律得

$p_{1}V_{1}=p_{2}V_{2}$，解得 $V_{2}=p_{1}V_{1}/p_{2}=1\times10^{-19}\Umc$，

则 $N=\frac{V_{2}}{V_{\text{标}}}N_{\text{A}}=\frac{1\times10^{-19}}{22.4\times10^{-3}}\times6.02\times10^{23}$ 个 $=3\times10^{6}$ 个。
}

\item
%% number: 23
%% paperName: 1988年全国高考第 23 题 3 分
%% typeId: 3
%% type: 填空题
%% chapter: 力物体的平衡
%% point: 力矩平衡
%% method: 无
%% score: 3
%% degree: 550
%% duplicateId: 0
%% body:
一均匀木杆，每米重 $10\UNm$，支点位于离木杆的左端点 $0.3\Um$ 处。现将一重量为 $11\UN$ 的物体挂在木杆的左端点上。设在木杆的右端点施一大小为 $5.0\UN$ 的竖直向上的力，恰能使木杆平衡，则木杆的长度 $L=$\tkanswer[1.3cm]{$1.8\Um$}。

%% answer: 1.8 m
\jdanswer{
$1.8\Um$
}

%% memo:
\memoanswer{
设木杆的长度为 $l$，由力矩平衡得 $11\times0.3+5(l-0.3)=10l\left(\frac{l}{2}-0.3\right)$，即 $5l^{2}-8l-1.8=0$，解得 $l=1.8\Um$。
}

\item
%% number: 24
%% paperName: 1988年全国高考第 24 题 3 分
%% typeId: 3
%% type: 填空题
%% chapter: 机械波
%% point: 波的多解性
%% method: 无
%% score: 3
%% degree: 750
%% duplicateId: 0
%% body:
绳上有一简谐横波向右传播，当绳上某质点 A 向上运动到最大位移时，在其右方相距 $0.30\Um$ 的质点 B 刚好向下运动到最大位移。已知波长大于 $0.15\Um$，则该波的波长等于多少米？答：\tkanswer[2.4cm]{$0.6\Um$，$0.2\Um$}。

%% answer: 0.6 m，0.2 m
\jdanswer{
$0.6\Um$，$0.2\Um$
}

%% memo:
\memoanswer{
满足要求 A、B 间的最短间距为 $\frac{\lambda}{2}$，由波的周期性知 $n\lambda+\frac{1}{2}\lambda=0.3$，解得 $\lambda=\frac{0.3}{n+\frac{1}{2}}$，当 $n=0$，$\lambda=0.6\Um$；当 $n=1$，$\lambda=0.2\Um$；当 $n=2$，$\lambda=0.12\Um$，因 $\lambda>0.15\Um$，所以波长为 $0.6\Um$ 或 $0.2\Um$。
}

\item
%% number: 25
%% paperName: 1988年全国高考第 25 题 3 分
%% typeId: 3
%% type: 填空题
%% chapter: 光学
%% point: 光的折射
%% method: 无
%% score: 3
%% degree: 550
%% duplicateId: 0
%% body:
在厚度为 $d$、折射率为 $n$ 的大玻璃板的下表面，紧贴着一个半径为 $r$ 的圆形发光面。为了从玻璃板的上方看不见圆形发光面，可在玻璃板的上表面贴一块纸片，所贴纸片的最小面积为\tkanswer[2.6cm]{$\pi\left(r+\frac{d}{\sqrt{n^{2}-1}}\right)^{2}$}。
\onepicture[align=right]{figs/25.png}

%% answer: \pi(r+\frac{d}{\sqrt{n^2-1}})^2
\jdanswer{
$\pi\left(r+\frac{d}{\sqrt{n^{2}-1}}\right)^{2}$
}

%% memo:
\memoanswer{
设发生全反射临界角为 $C$，则 $\sin C=\frac{1}{n}$，临界情况为发光面边缘发出的光，经过玻璃板向空气中折射时，折射角为 $90^{\circ}$，此时的入射角为 $C$，设发光点与折射点的水平间距为 $x$，由几何关系得 $\tan C=\frac{x}{d}$，解得 $x=d\tan C=\frac{d}{\sqrt{n^{2}-1}}$，又 $R=r+x$，故 $S=\pi R^{2}=\pi\left(r+\frac{d}{\sqrt{n^{2}-1}}\right)^{2}$。
}

\item
%% number: 26
%% paperName: 1988年全国高考第 26 题 10 分
%% typeId: 4
%% type: 实验题
%% chapter: 电路
%% point: 电路实验
%% method: 无
%% score: 10
%% degree: 450
%% duplicateId: 0
%% body:
右图是测定电流表内电阻实验的电路图。电流表的内电阻约在 $100\UO$ 左右，满偏电流为 $500\UuA$。用电池作电源。

\begin{enumerate}
	\item
	实验室中配有的可变电阻为：

	（A）电阻箱，阻值范围为 $0\sim10\UO$。

	（B）电阻箱，阻值范围为 $0\sim9999\UO$。

	（C）电位器，阻值范围为 $0\sim200\UO$。

	（D）电位器，阻值范围为 $0\sim20\UkO$。

	在上述配有的可变电阻中，电路图中的 $R$ 应选用\tkanswer[0.8cm]{D}，$R'$ 应选用\tkanswer[0.8cm]{B}。（填写字母代号）
	\item
	某学生进行的实验步骤如下：

	\begin{steps}
		\item
		先将 $R$ 的阻值调到最大，合上 $K_{1}$，调节 $R$ 的阻值，使电流表的指针偏转到满刻度。
		\item
		合上 $K_{2}$，调节 $R'$ 和 $R$ 的阻值，使电流表的指针偏转到满刻度的一半。
		\item
		记下 $R'$ 的阻值。
	\end{steps}

	指出上述实验步骤中有什么错误。答：\tkanswer[6cm]{合上 $K_{2}$ 后，不应再调节 $R$ 的阻值}。
	\item
	如果按正确实验步骤测得的 $R'$ 值为 $100\UO$，已知电流表的满偏电流为 $500\UuA$，现在要把它改装成量程为 $2\UV$ 的伏特表，则应串联的分压电阻为\tkanswer[1.4cm]{$3900\UO$}。
	\item
	由以下四种说法中选出正确说法，并将正确说法前的字母填在题后方括号内。

	（A）上述电流表内电阻的测得值比其真实值小，这一因素使得用上述改装成的伏特表测定电压时，其读数比伏特表两端的实际电压小

	（B）上述电流表内电阻的测得值比其真实值大，这一因素使得用上述改装成的伏特表测定电压时，其读数比伏特表两端的实际电压小

	（C）上述电流表内电阻的测得值比其真实值小，这一因素使得用上述改装成的伏特表测定电压时，其读数比伏特表两端的实际电压大

	（D）上述电流表内电阻的测得值比其真实值大，这一因素使得用上述改装成的伏特表测定电压时，其读数比伏特表两端的实际电压大

	答：\tkanswer[0.8cm]{A}。
\end{enumerate}
\onepicture[align=right]{figs/26.png}

%% answer: （1）D，B；（2）合上 K2 后，不应再调节 R 的阻值；（3）3900；（4）A
\jdanswer{
\begin{enumerate}
	\item
	D，B

	\item
	合上 $K_{2}$ 后，不应再调节 $R$ 的阻值

	\item
	3900

	\item
	A
\end{enumerate}
}

%% memo:
\memoanswer{
（1）因要读出电阻的阻值，所以应使用电阻箱，因 $R_{g}\approx100\UO$，A 项的阻值不够用，故 B 正确；因电阻 $R=\frac{E}{I}=\frac{3}{500\times10^{-6}}\UkO\approx6\UkO$；C 项的阻值不够用，故 D 正确。

（2）因电流表半偏时认为是满偏时电流的一半，所以②中不应再调节 $R$。

（3）电流表的满偏电压为 $U_{\mathrm{g}}=I_{\mathrm{g}}R_{\mathrm{g}}=500\times10^{-6}\times100\UV=0.05\UV$，故与之串联的分压电阻为 $\frac{R}{R_{g}}=\frac{U-U_{g}}{U_{g}}=\frac{2}{0.05}-1=39$，解得 $R=3900\UO$。

（4）因为合上 $K_{2}$ 后，$R_{g}$ 并联一个 $R'$，总电流比满偏电流大，所以 $R'$ 上电流大于半偏，并联电路电阻与电流成反比，所以 $R'<R_{g}$，故 BD 错误；因此测得值小于真实值，即 $R'<R_{g}$，所以 $I(R'+R)<I(R_{g}+R)$，即读数小于真实值。例如：$R'=100\UO$，$R_{g}=110\UO$，当 $I_{g}=500\UuA$ 时，改装的电压表读数为 2 V，实际电压为 2.005 V，故 A 正确 B 错误。
}

\item
%% number: 27
%% paperName: 1988年全国高考第 27 题 3 分
%% typeId: 4
%% type: 实验题
%% chapter: 力物体的平衡
%% point: 力
%% method: 无
%% score: 3
%% degree: 650
%% duplicateId: 0
%% body:
有一架托盘天平，没有游码，最小砝码为 $100\Umg$。用这架天平称量一个物体，当右盘中加上 $36.20\Ug$ 砝码时，天平指针向左偏 1.0 小格，如图中实箭头所示。如果在右盘中再加上 $100\Umg$ 的砝码，天平指针则向右偏 1.5 小格，如图中虚箭头所示。这个物体的质量可读为\tkanswer[1.5cm]{$36.24\Ug$}。
\onepicture[align=right]{figs/27.png}

%% answer: 36.24 g
\jdanswer{
$36.24\Ug$
}

%% memo:
\memoanswer{
$100\Umg$ 对应 2.5 格，每格对应 $40\Umg$，$m=36.20\Ug+0.040\Ug=36.24\Ug$。
}

\item
%% number: 28
%% paperName: 1988年全国高考第 28 题 7 分
%% typeId: 6
%% type: 计算题
%% chapter: 运动和力
%% point: 动量和能量
%% method: 无
%% score: 7
%% degree: 650
%% duplicateId: 0
%% body:
用细线悬挂一质量为 $M$ 的木块，木块静止，如下左图所示。现有一质量为 $m$ 的子弹自左方水平地射穿此木块，穿透前后子弹的速度分别为 $v_{0}$ 和 $v$。求木块能摆到的最大高度。（设子弹穿过木块的时间很短，可不计）
\onepicture[align=right]{figs/28.png}

%% answer: h=\frac{1}{2g}\left(\frac{mv_0-mv}{M}\right)^2
\jdanswer{
$h=\frac{1}{2g}\left(\frac{mv_{0}-mv}{M}\right)^{2}$
}

%% memo:
\memoanswer{
射穿过程中，水平方向动量守恒，可得

$mv_{0}=MV+mv$ ①

射穿后，木块在摆动过程中机械能守恒，可得

$\frac{1}{2}Mv^{2}=Mgh$ ②

由以上两式可得

$h=\frac{1}{2g}\left(\frac{mv_{0}-mv}{M}\right)^{2}$ ③

评分标准：本题共 7 分。列出（1）式给 3 分；列出（2）式给 3 分；得出（3）式给 1 分。
}

\item
%% number: 29
%% paperName: 1988年全国高考第 29 题 9 分
%% typeId: 6
%% type: 计算题
%% chapter: 气体性质
%% point: 玻意耳定律
%% method: 无
%% score: 9
%% degree: 550
%% duplicateId: 0
%% body:
一圆筒形气缸静置于地面上，如右图所示。气缸筒的质量 $M$，活塞（连同手柄）的质量为 $m$，气缸内部的横截面积为 $S$。大气压强为 $p_{0}$。平衡时气缸内的容积为 $V$。现用手握住活塞手柄缓慢向上提。设气缸足够长，在整个上提过程中气体温度保持不变，并不计气缸内气体的重量及活塞与气缸壁间的摩擦。求将气缸刚提离地面时活塞上升的距离。
\onepicture[align=right]{figs/29.png}

%% answer: x=\frac{(M+m)gV}{(p_0S-Mg)S}
\jdanswer{
$x=\frac{(M+m)gV}{(p_{0}S-Mg)S}$
}

%% memo:
\memoanswer{
设气缸内气体原来的压强为 $p_{1}$，后来的压强为 $p_{2}$，由活塞的受力平衡得

$p_{1}S=p_{0}S+mg$ ①

$p_{0}S=p_{2}S+Mg$ ②

由玻意耳-马略特定律可知

$p_{1}V=p_{2}(V+xS)$ ③

其中 $x$ 为活塞上升的距离。由（1）（2）（3）式可得

$x=\frac{(M+m)gV}{(p_{0}S-Mg)S}$ ④

评分标准：本题共 9 分。列出（1）式给 2 分；列出（2）式给 2 分；列出（3）式给 3 分；得出（4）式再给 2 分。
}

\item
%% number: 30
%% paperName: 1988年全国高考第 30 题 9 分
%% typeId: 6
%% type: 计算题
%% chapter: 电场
%% point: 带电粒子的直线运动
%% method: 无
%% score: 9
%% degree: 350
%% duplicateId: 0
%% body:
$N$ 个长度逐个增大的金属圆筒和一个靶，它们沿轴线排列成一串，如图所示（图中只画出了六个圆筒，作为示意）。各筒和靶相间地连接到频率为 $\nu$、最大电压值为 $U$ 的正弦交流电源的两端。整个装置放在高真空容器中。圆筒的两底面中心开有小孔。现有一电量为 $q$、质量为 $m$ 的正离子沿轴线射入圆筒，并将在圆筒间及圆筒与靶间的缝隙处受到电场力的作用而加速（设圆筒内部没有电场）。缝隙的宽度很小，离子穿过缝隙的时间可以不计。已知离子进入第一个圆筒左端的速度为 $v_{1}$，且此时第一、二两个圆筒间的电势差 $U_{1}-U_{2}=-U$。为使打到靶上的离子获得最大能量，各个圆筒的长度应满足什么条件？并求出在这种情况下打到靶上的离子的能量。
\onepicture[align=right]{figs/30.png}

%% answer: L_n=\frac{1}{2\nu}\sqrt{\frac{2(n-1)qU}{m}+v_1^2}；E_k=NqU+\frac{1}{2}mv_1^2
\jdanswer{
\begin{enumerate}
	\item
	$L_{n}=\frac{1}{2\nu}\sqrt{\frac{2(n-1)qU}{m}+v_1^{2}}$（$n=1,2,3\ldots\ldots N$）

	\item
	$E_{k}=NqU+\frac{1}{2}mv_{1}^{2}$
\end{enumerate}
}

%% memo:
\memoanswer{
为使正离子获得最大能量，要求离子每次穿越缝隙时，前一个圆筒的电势比后一个圆筒的电势高 $U$，这就要求离子穿过每个圆筒的时间都恰好等于交流电的半个周期。由于圆筒内无电场，离子在筒内做匀速运动。设 $v_{n}$ 为离子在第 $n$ 个圆筒内的速度，则有

$qU=\frac{1}{2}mv_{n+1}^{2}-\frac{1}{2}mv_{n}^{2}$ ①

第 $n$ 个圆筒的长度为 $L_{n}=v_{n}\frac{T}{2}=\frac{v_{n}}{2\nu}$ ②

由（1）式得 $(n-1)qU=\frac{1}{2}mv_{n}^{2}-\frac{1}{2}mv_{1}^{2}$，解得 $v_{n}=\sqrt{\frac{2(n-1)qU}{m}+v_{1}^{2}}$ ③

将（3）代入（2），得第 $n$ 个圆筒的长度应满足的条件为：

$L_{n}=\frac{1}{2\nu}\sqrt{\frac{2(n-1)qU}{m}+v_{1}^{2}}$ ④

$n=1,2,3\ldots\ldots N$。

打到靶上的离子的能量为：

$E_{k}=NqU+\frac{1}{2}mv_{1}^{2}$ ⑤

评分标准：本题共 9 分。列出（1）式给 2 分；列出（2）式给 3 分；得出（4）式再给 2 分；得出（5）式给 2 分。
}

\end{enumerate}

\end{document}
